\documentclass[10pt,a4paper]{amsart}
\usepackage[margin=19mm]{geometry}
\usepackage{amsmath,amssymb,mathtools}
\usepackage[scr=rsfs,cal=euler]{mathalfa}
\usepackage{xfrac}
\usepackage[unicode,hidelinks,bookmarks=false]{hyperref}
\newcommand{\ie}{i.e.\ }
\newcommand{\cf}{cf.\ }
\newcommand{\resp}{respectively\ }
\newcommand{\Subsection}{\subsection}
\providecommand{\nobreakdash}{\nobreak\hbox{-}}
\setlength{\emergencystretch}{2em}
\multlinegap0pt
\usepackage{dsfont}
\usepackage{mathtools}
\usepackage{amssymb}
\usepackage{xcolor}
\usepackage{tensor}
\usepackage[normalem]{ulem}
\newtheorem{proposition}{Proposition}
\title{Detecting screens modeled by Schr\"odinger operators that generate $C_0$ contraction semigroups}
\author{Lawrence Frolov}
\date{Source: arXiv:2506.00231v2 (CC BY 4.0)}
\begin{document}
\maketitle

\noindent\emph{Source and licence.} Detecting screens modeled by Schr\"odinger operators that generate $C_0$ contraction semigroups. Version arXiv:2506.00231v2 (CC BY 4.0), distributed by arXiv under the Creative Commons Attribution 4.0 International licence, http://creativecommons.org/licenses/by/4.0/, as recorded in the arXiv licence metadata for this exact version and shown on the abstract page https://arxiv.org/abs/2506.00231v2. Full licence notice: \texttt{licenses/arxiv-2506.00231-CC-BY-4.0.txt}.\par\smallskip

\noindent\textit{Editorial scope.} Counted unit: the article's proposition Decomposition (source0.tex lines 518--530). The article's complete original proof of it is delivered with it (source0.tex lines 531--564, one \texttt{proof} environment of 34 lines). No mathematics is written, repaired or re-derived: the statement and the proof are the article's own lines, in the article's own notation, and no retained line is substituted or re-flowed.  Retained uncounted as the source-local context the proof needs: lines 488--490; lines 512--514.  Nothing was omitted from any retained span, so the package carries no omission record.  No retained line contains a citation, so the wrapper carries no bibliography and no bibliography entry.  No retained line contains a figure environment, an image-inclusion command or drawing code, so no image is delivered and the package declares no asset dependency.  Editorial formatting only, in addition: an AMS \texttt{amsart} preamble that restates the article's own declarations and macros used by the retained lines, namely the environments \texttt{proposition} on separate counters, together with the standard packages those lines need.  The retained environments print in this document as Proposition 1 in order of appearance, whereas the article prints the same items with the numbering of its own full document.  The complete native source of the article as distributed in the arXiv e-print for this exact version is bundled unchanged as \texttt{sources/179069/source0.tex}.
    \begin{equation}\label{Dirichlet Star Decomp}
        D(\hat{H}^*)= \ker \tau_D \oplus \ker(\hat{H}^*-\lambda)
    \end{equation}

    \begin{equation}\label{Green's Identity Star}
        \langle -i\hat{H}^*\psi,v\rangle_{L^2(\Omega)}+ \langle \psi,-i\hat{H}^*v \rangle_{L^2(\Omega)}=i\left( \langle \tilde{\tau}_N \psi, \tau_D v\rangle_{H^{-3/2}(\partial \Omega)\times H^{3/2}(\partial \Omega)} -\langle  \tilde{\tau}_D \psi, \tau_N v\rangle_{H^{-1/2}(\partial \Omega)\times H^{1/2}(\partial \Omega)}\right).
    \end{equation}

\begin{proposition}\label{Current Decomposition}
        Let $\eta \in \rho(\hat{H}_D)\cap \mathbb{R}$, which exists since $\hat{H}_D-V$ is a positive operator and $V \in L^\infty(\Omega,\mathbb{R})$. Also let $\iota_\pm:H^{\pm {1/2}}(\partial \Omega)\to L^2(\partial \Omega)$ be the isomorphisms such that $\langle \xi,\chi \rangle_{H^{-1/2}(\partial \Omega) \times H^{1/2}(\partial \Omega)}=\langle \iota _- \xi, \iota_+ \chi \rangle_{L^2(\partial \Omega)}$. Then, defining $G_\pm(\eta): D(-i\hat{H}^*)\to L^2(\partial \Omega)$ by
    \begin{equation}
        G_+\psi:=\frac{1}{\sqrt{2}}\left(\iota_- \tilde{\tau}_D \psi  + i \iota_+ \tau_N \psi_D\right)
    \end{equation}
    \begin{equation}
        G_-\psi:=\frac{1}{\sqrt{2}}\left(\iota_-\tilde{\tau}_D \psi  - i \iota_+ \tau_N  \psi_D \right)
    \end{equation}
    where $\psi=\psi_D + \psi_\eta$ is decomposed according to equation (\ref{Dirichlet Star Decomp}) returns a further extension of the Green's identity for elements in $\psi, \phi \in D(\hat{H}^*)$
    \begin{equation}\label{Current Star Decomp Eq}
        \langle -i\hat{H}^* \psi,\phi \rangle_{L^2(\Omega)} + \langle \psi,-i\hat{H}^* \phi \rangle_{L^2(\Omega)}=\langle G_+\psi,G_+\phi \rangle_{L^2(\partial \Omega)} - \langle G_-\psi,G_- \phi \rangle_{L^2(\partial \Omega)}.
    \end{equation}
        \end{proposition}

      \begin{proof}
          Let $\psi,\phi \in D(\hat{H})$, and decompose $\psi= \psi_D + \psi_\eta$, $\phi=\phi_D + \phi_\eta$ according to equation (\ref{Dirichlet Star Decomp}). First, since $\hat{H}_D$ is self-adjoint 
          \begin{equation}
           \langle -i\hat{H}^* \psi_D,\phi_D \rangle_{L^2(\Omega)} + \langle \psi_D,-i\hat{H}^* \phi_D \rangle_{L^2(\Omega)}=\langle -i\hat{H}_D \psi_D,\phi_D \rangle_{L^2(\Omega)} + \langle \psi_D,-i\hat{H}_D \phi_D \rangle_{L^2(\Omega)}=0.
          \end{equation}
          Similarly, since $\eta$ is real we have 
          \begin{equation}
           \langle -i\hat{H}^* \psi_\eta,\phi_\eta \rangle_{L^2(\Omega)} + \langle \psi_\eta,-i\hat{H}^* \phi_\eta \rangle_{L^2(\Omega)}=\langle -i\eta\psi_\eta,\phi_\eta \rangle_{L^2(\Omega)} + \langle \psi_\eta,-i\eta \phi_\eta \rangle_{L^2(\Omega)}=0.
           \end{equation}
           Hence
\begin{equation}
    \begin{split}
             \langle -i\hat{H}^* \psi,\phi \rangle_{L^2(\Omega)} + \langle \psi,-i\hat{H}^* \phi \rangle_{L^2(\Omega)}&=\langle -i\hat{H}^* \psi_\eta,\phi_D \rangle_{L^2(\Omega)} +\langle \psi_\eta,-i\hat{H}^* \phi_D \rangle_{L^2(\Omega)}
             \\
             &+ \langle -i\hat{H}^* \psi_D,\phi_\eta \rangle_{L^2(\Omega)} +\langle \psi_D,-i\hat{H}^* \phi_\eta \rangle_{L^2(\Omega)}.
    \end{split}
\end{equation}
Since $\psi_D, \phi_D \in H^2(\Omega)$, we may apply the generalized Green's identity $(\ref{Green's Identity Star})$ to the two pairs of terms. Applying the identity along with $\tau_D \psi_D= \tau_D \phi_D=0$ returns
\begin{equation}\label{Gen Green's Identity Inner Prod}
\begin{split}
    \langle -i\hat{H}^* \psi,\phi \rangle_{L^2(\Omega)} + \langle \psi,-&i\hat{H}^* \phi \rangle_{L^2(\Omega)}=
    \\
    &=i\left( \langle \tau_N \psi_D, \tilde{\tau}_D \phi_\eta \rangle_{H^{1/2}(\partial \Omega)\times H^{-1/2}(\partial \Omega)} -\langle  \tilde{\tau}_D \psi_\eta, \tau_N \phi_D \rangle_{H^{-1/2}(\partial \Omega)\times H^{1/2}(\partial \Omega)}\right)
    \\
    &=i\left( \langle \tau_N \psi_D, \tilde{\tau}_D \phi \rangle_{H^{1/2}(\partial \Omega)\times H^{-1/2}(\partial \Omega)} -\langle  \tilde{\tau}_D \psi, \tau_N \phi_D \rangle_{H^{-1/2}(\partial \Omega)\times H^{1/2}(\partial \Omega)}\right)
    \\
    &=i\left( \langle \iota_+ \tau_N \psi_D, \iota_- \tilde{\tau}_D \phi \rangle_{L^2(\partial \Omega)} -\langle  \iota_-\tilde{\tau}_D \psi, \iota_+ \tau_N \phi_D \rangle_{L^2(\partial \Omega)}\right).
    \end{split}
\end{equation}
It is then not difficult to compute
    \begin{equation}
        \langle G_+\psi, G_+ \phi\rangle_{L^2(\partial \Omega)}- \langle G_-\psi, G_- \phi\rangle_{L^2(\partial \Omega)}= i\left( \langle \iota_+ \tau_N \psi_D, \iota_- \tilde{\tau}_D \phi \rangle_{L^2(\partial \Omega)} -\langle  \iota_-\tilde{\tau}_D \psi, \iota_+ \tau_N \phi_D \rangle_{L^2(\partial \Omega)}\right).
    \end{equation}
This concludes our proof of Proposition \ref{Current Decomposition}. \end{proof}

\end{document}
